% !TeX root = ../icml2026.tex
\section{Simulation Code}\label{app:code}

\begin{lstlisting}
import torch
from torch import optim
from tqdm import trange
import itertools
import matplotlib.pyplot as plt
import numpy as np

class Trainer:
    def __init__(self, n, k):
        self.n = n
        self.k = k
        self.device="cuda" if torch.cuda.is_available() else "cpu"
        print("Using device", self.device)
        self.all_combinations = list(itertools.combinations(range(self.n), self.k))
        self.subset_indices_tensor = torch.tensor([list(s) for s in self.all_combinations], device=self.device)
        self.subset_excluded_tensor = torch.tensor(
            [
                [i for i in range(self.n) if i not in subset_indices]
                for subset_indices in self.all_combinations
            ], device=self.device
        )
        self.total_violations = len(self.all_combinations) * self.k * (self.n-self.k)

    def calculate_loss(self):
        # Convert batch_subset_indices to a tensor for efficient indexing
        # Calculate sums of vectors for each subset in the batch
        # Shape: (batch_size, d)
        subset_sums = self.vector_embeddings[self.subset_indices_tensor].mean(dim=1)
        # Calculate dot products of subset sums with all vectors
        # Shape: (batch_size, n)
        dot_products_all = torch.matmul(subset_sums, self.vector_embeddings.T)
        # Get dot products with vectors in the subset by indexing dot_products_all
        # Shape: (batch_size, k)
        dot_products_xi = torch.gather(dot_products_all, 1, self.subset_indices_tensor)
        # For each subset in the batch, calculate dot products with vectors NOT in the subset
        # We can create a mask to zero out elements within the subset
        # This is faster than iterating through the batch
        # Shape: (batch_size, n-k)
        dot_products_xj = torch.gather(dot_products_all, 1, self.subset_excluded_tensor)
        # Calculate differences and apply ReLU
        # We need to unsqueeze dot_products_xi to match dimensions for broadcasting
        # Shape: (batch_size, n-k) - (batch_size, k)
        differences = dot_products_xi.unsqueeze(1) - dot_products_xj.unsqueeze(2)
        total_loss = torch.relu(-differences).sum((-1, -2)).mean()
        num_violations = (differences < 0).sum().item()
        return total_loss, num_violations


    def train(self, d, num_epochs, learning_rate=0.1, patience=20):
        self.vector_embeddings = torch.randn(
            self.n, d, device=self.device, requires_grad=True
        )


        optimizer = optim.Adam(
            [self.vector_embeddings],
            lr=learning_rate,
        )
        scheduler = optim.lr_scheduler.OneCycleLR(
            optimizer=optimizer,
            max_lr=learning_rate,
            total_steps=num_epochs,
            pct_start=0.0,
        )

        min_violations = self.total_violations
        epochs_no_improve = 0

        with trange(num_epochs, desc=f"\t\t n={self.n}, k={self.k}, {d=}") as epoch_iterator:
            for epoch in epoch_iterator:
                optimizer.zero_grad()

                loss, violations = self.calculate_loss()
                loss.backward()
                optimizer.step()
                scheduler.step()


                epoch_loss = loss.item()
                epoch_iterator.set_postfix(
                    {
                        "n": self.n,
                        "k": self.k,
                        "loss": epoch_loss,
                        "#vio rate": violations / self.total_violations,
                        "min #vio": min_violations,
                        "epochs_no_improve": epochs_no_improve
                    }
                )

                if violations < min_violations:
                    min_violations = violations
                    epochs_no_improve = 0
                else:
                    epochs_no_improve += 1

                if violations == 0:
                    print("Early stopping: No violations found.")
                    break
                if epochs_no_improve >= patience:
                    print(f"Early stopping: No improvement in violations for {patience} epochs.")
                    break

        return min_violations


class Experiment:
    def __init__(self):
        self.search_paths = {}
        self.minimal_dimensions = {}

    def find_minimal_dimension(self, k, n_values, left0=0, num_epochs=100, learning_rate=0.1, patience=10):

        last_minimal = left0
        for n in n_values:
            print("#" * 10 + "new task" + "#" * 10)
            print(f"Finding minimal dimension for n={n}, k={k}")
            print("#" * 30)
            trainer = Trainer(n, k)

            self.search_paths[n] = []
            left, right = last_minimal+1, last_minimal + 40 # the range is set by prior to accelerate the convergence.
            minimal_d = n + 1

            while left <= right:
                mid = (left + right) // 2
                if mid == 0:
                    mid = 1

                print(f"\t>>>>Testing dimension d={mid}")

                violations = trainer.train(
                    d=mid,
                    num_epochs=num_epochs,
                    learning_rate=learning_rate / np.log2(n),
                    patience=patience
                )

                print(f"\t<<<<Violations for d={mid}: {violations}")

                self.search_paths[n].append({'dimension': mid, 'violations': violations})

                if violations == 0:
                    minimal_d = mid
                    right = mid - 1
                else:
                    left = mid + 1


            self.minimal_dimensions[n] = minimal_d if minimal_d <= n else -1
            last_minimal = minimal_d

            with open("minimal_dem_log.txt", "at") as f:
                f.write(f"minimal dimension @ {k=}&{n=} is {minimal_d}\n")
            print("#" * 10 + " Task Finished" + "#" * 10)
            print(f"minimal dimension @ {k=}&{n=} is {minimal_d}\n")
            print("#" * 30)

        return self.minimal_dimensions

    def plot_minimal_dimension_vs_n(self, k):
        if not self.minimal_dimensions:
            print("No experiment results to plot. Run find_minimal_dimension first.")
            return

        n_plot = list(self.minimal_dimensions.keys())
        d_plot = list(self.minimal_dimensions.values())

        plt.figure(figsize=(8, 6))
        plt.plot(n_plot, d_plot, marker='o', linestyle='-')
        plt.xlabel("Number of Vectors (n)")
        plt.ylabel("Minimal Dimension (d)")
        plt.xscale('log')
        plt.title(f"Minimal Dimension vs. Number of Vectors for k={k}")
        plt.grid(True)
        plt.show()

        plt.savefig("med_vs_n.png")

    def plot_violations_vs_d(self, k):
        if not self.search_paths:
            print("No experiment results to plot. Run find_minimal_dimension first.")
            return

        for n, path in self.search_paths.items():
            sorted_path = sorted(path, key=lambda x: x['dimension'])

            dimensions = [item['dimension'] for item in sorted_path]
            violations = [item['violations'] for item in sorted_path]

            plt.figure(figsize=(8, 6))
            plt.plot(dimensions, violations, marker='o', linestyle='-')
            plt.xlabel("Dimension (d)")
            plt.ylabel("Number of Violations")
            plt.title(f"Violations vs. Dimension for n={n}, k={k}")
            plt.grid(True)
            plt.show()


# 1. Define a list of n_values to experiment with
n_values = [5 * 2 ** (i) for i in [0, 1, 2, 3, 4, 5]]

# 2. Set a value for k
k = 2

# 3. Instantiate the Experiment class
experiment = Experiment()

# 4. Call the find_minimal_dimension method with a reasonable batch_size and limited epochs
#    Using a batch_size of 32 and num_epochs of 100 for initial testing.
print("Starting experiment with mini-batching and early stopping...")
minimal_dims = experiment.find_minimal_dimension(
    k, n_values, learning_rate=1, num_epochs=1000, patience=1000)

# 5. Print the minimal_dims dictionary
print("\n Minimal dimensions found:", minimal_dims)

# 6. Call the plot_minimal_dimension_vs_n method
experiment.plot_minimal_dimension_vs_n(k)

# 7. Call the plot_violations_vs_d method
experiment.plot_violations_vs_d(k)
\end{lstlisting}
